\section{研究计划进度安排及预期目标}

\subsection{进度安排}

3.17 – 3.24 完成预处理中的 LOD 生成部分，并针对不同场景进行测试；

3.24 – 3.31 利用网格着色器实现剔除和 LOD 选择，并进行测试；

4.1  -  4.15 实现基本的流式加载；

4.16 -  4.23 对流式加载的性能进行优化，保证显存占用和帧率的平衡；

4.24 -  4.30 优化生成 LOD 的速度和效果；

5.1  -  5.8  实现异步加载；

5.9  -  5.16 针对特定工业场景进行全流程的优化；

5.17 -  5.24 完成全流程，在不同场景下与虚幻引擎的 Nanite 进行对比测试，分析设计的优点与不足之处，提出未来改进方案。

\subsection{预期目标}

（1）支持包含数亿级多边形的工业模型导入，并确保数据预处理与加载过程在 1 小时内完成，同时内存占用不超过 24GB ；

（2）在运行过程中，显存占用不超过 8GB ，并确保渲染帧率稳定在 30 FPS 以上，同时保持较高的画面质量，使渲染结果与原始模型的视觉差异控制在可接受范围内（PSNR 不低于一定阈值）；

（3）在帧率、显存占用、内存占用、预处理时间以及画面保真度等关键指标上，至少达到与虚幻引擎 Nanite 相当的水平，并在部分方面实现超越；

（4）具有较强的鲁棒性，能够适应不同类型的工业场景及 CAD 模型，并可根据不同硬件配置进行优化。